% Fragment do pracy: \input{figures/uklad_pomiarowy/uklad_pomiarowy_body.tex}
% Wymaga w preambule: \usepackage[european,siunitx]{circuitikz}
% do pliku figures/... i wstaw przez \input w środowisku figure.
\documentclass[border=4mm]{standalone}
\usepackage[T1]{fontenc}
\usepackage[utf8]{inputenc}
\usepackage[polish]{babel}
\usepackage{lmodern}
\usepackage[european,siunitx]{circuitikz}
\usepackage{siunitx}
\sisetup{output-decimal-marker={,}}
\begin{document}
%%% START
\begin{circuitikz}[scale=1, transform shape, line width=0.6pt]
  \ctikzset{bipoles/length=1.2cm}
  % --- Zasilacz 1 (obwód targetu) ---
  \draw (0,0) to[battery1, invert, l_={$U_\mathrm{T}$}] (0,6);
  \node[left, align=right, font=\small] at (-0.6,3) {zasilacz 1\\(target)};
  \node[left] at (-0.25,5.4) {$+$};
  \node[left] at (-0.25,0.6) {$-$};
  % rezystor ograniczający z woltomierzem równolegle
  \draw (0,6) -- (1.2,6) to[R, l={$R=\SI{100}{\kilo\ohm}$}, *-*] (4.2,6) -- (5.4,6);
  % woltomierz mierzący napięcie na R
  \draw (1.2,6) -- (1.2,4.6) to[rmeter, t=V, l_={$U_R$}] (4.2,4.6) -- (4.2,6);
  \draw (5.4,6) -- (8,6) -- (8,5.15);

  % --- Źródło RTG MEMS (komora próżniowa) ---
  \draw[dashed, rounded corners=3pt] (6.6,1.0) rectangle (9.4,5.6);
  \node[anchor=north west, font=\scriptsize, align=left] at (6.65,5.55) {próżnia};
  % target (T)
  \draw[line width=2.4pt] (7.2,5.0) -- (8.8,5.0);
  \node[right] at (8.85,5.0) {T};
  % bramka (B) z aperturą
  \draw[line width=1.6pt] (7.0,3.3) -- (7.8,3.3);
  \draw[line width=1.6pt] (8.2,3.3) -- (9.0,3.3);
  \node[above] at (8.75,3.35) {B};
  % emiter (E) - ostrze
  \draw[line width=1.6pt] (7.3,1.6) -- (8.7,1.6);
  \filldraw (7.75,1.6) -- (8,2.25) -- (8.25,1.6) -- cycle;
  \node[left] at (7.3,1.6) {E};
  % wiązka elektronów
  \draw[-latex, densely dotted] (8,2.35) -- (8,4.85);
  \node[font=\scriptsize, rotate=90] at (7.75,4.1) {$e^-$};

  % --- Zasilacz 2 (bramka, B_on) ---
  \draw (9.0,3.3) -- (10.6,3.3) -- (10.6,6) -- (12,6);
  \draw (12,0) to[battery1, invert, l_={$U_\mathrm{B}$}] (12,6);
  \node[right, align=left, font=\small] at (12.9,3) {zasilacz 2\\(bramka, $B_\mathrm{on}$)};
  \node[right] at (12.25,5.4) {$+$};
  \node[right] at (12.25,0.6) {$-$};

  % --- Wspólna szyna (emiter) ---
  \draw (8,1.6) -- (8,0) node[circ]{};
  \draw (0,0) -- (8,0);
  % gałąź powrotna zasilacza 2: rezystor R i amperomierz
  \draw (8,0) -- (8.6,0) to[R, l={$R_\mathrm{B}=\SI{100}{\kilo\ohm}$}, *-*] (10.4,0)
        -- (12,0);
  % woltomierz mierzący napięcie na R_B
  \draw (8.6,0) -- (8.6,-1.4) to[rmeter, t=V, l_={$U_{R_\mathrm{B}}$}] (10.4,-1.4) -- (10.4,0);
\end{circuitikz}
